% Generated draft: requires mathematical and source review.
\documentclass[11pt]{article}
\usepackage[margin=25mm]{geometry}
\usepackage{fontspec}
\setmainfont{Latin Modern Roman}
\usepackage{amsmath,amssymb,mathtools,unicode-math}
\setmathfont{Latin Modern Math}
\usepackage{enumitem,needspace,xurl,hyperref}
\setlength{\parindent}{0pt}
\setlength{\parskip}{5pt}
\setlength{\emergencystretch}{3em}
\newcommand{\N}{\mathbb{N}}
\newcommand{\Z}{\mathbb{Z}}
\newcommand{\Q}{\mathbb{Q}}
\newcommand{\R}{\mathbb{R}}
\newcommand{\C}{\mathbb{C}}
\newcommand{\F}{\mathbb{F}}
\newcommand{\A}{\mathbb{A}}
\newcommand{\PP}{\mathbb{P}}
\newcommand{\E}{\mathbb{E}}
\newcommand{\eps}{\varepsilon}
\newcommand{\dd}{\,\mathrm{d}}
\begin{document}
\section*{UnsolvedMath ID: 30002102}
\textbf{Source problem code:} catalogue entry
\subsection{Definitions and mathematical statement}
Let \(k\) be an algebraically closed field and let \(X\) be a projective variety over \(k\). The record specifies neither a restriction on the characteristic of \(k\) nor a smoothness hypothesis on \(X\); its title explicitly concerns singular varieties.

\textbf{Requested construction.} Define a theory assigning to projective varieties \(X\) a graded ring denoted \(CH^*(X)\), with the expected contravariant functoriality, cycle-class and Chern-class properties, and a rational Chern-character isomorphism.

Under the usual interpretation of contravariant functoriality, a morphism \(f\colon X\to Y\) in the intended category would induce a graded ring homomorphism \(f^*\colon CH^*(Y)\to CH^*(X)\), satisfying
\[
(\operatorname{id}_X)^*=\operatorname{id}_{CH^*(X)},\qquad
(g\circ f)^*=f^*\circ g^*
\]
for composable morphisms \(f\colon X\to Y\) and \(g\colon Y\to Z\). The source does not specify the intended class of morphisms, so these identities describe the usual meaning of functoriality without resolving its scope.

\textbf{Unresolved specifications.} The word ``expected'' is not accompanied by a complete list of axioms. In particular, the record does not specify whether the underlying groups must be the usual groups of algebraic cycles modulo rational equivalence, the target and compatibility requirements of a cycle-class map, the objects admitting Chern classes, or the source, grading conventions, and compatibility requirements of the rational Chern-character isomorphism. These choices cannot be supplied silently for singular varieties.

The supplied background assessment instead describes an incomplete prompt about Chern classes and cycles supported in a given codimension, and reports that surrounding text exhibits failures of a naive formula without stating a replacement theorem. Neither that formula nor its hypotheses are supplied. Consequently, this is an incomplete construction problem: the record does not determine a fully specified existence assertion or establish how the codimension formulation relates to the stated ring-construction question.
\subsection{Short English statement}
For projective varieties over an algebraically closed field, including singular varieties, the record asks for a Chow-ring theory with pullbacks, cycle classes, Chern classes, and a rational Chern-character isomorphism. The required axioms and several domains are unspecified, and the supplied background describes a potentially different codimension question. A precise complete conjecture therefore requires further source clarification.
\subsection{Sources}
\begin{itemize}
\item Local UnsolvedMath registry, supplied complete source record: permanent ID \(30002102\), title ``Functorial Chow Rings for Singular Varieties,'' exact source problem code \texttt{catalogue entry}; collection: Oberwolfach Reports Open Problems. No individual proposer is supplied; the proposed year is recorded as 2012.
\item Supplied report citation: \emph{Classical Algebraic Geometry (2013)}, Oberwolfach Reports, DOI 10.4171/owr/2012/30, pp.\ 1845--1893. Supplied source URL and literature-source link, labelled ``original OWR report'': \url{https://doi.org/10.4171/owr/2012/30}. The citation's 2013 designation and the record's 2012 year information are retained as supplied.
\item Literature assessment dated exactly \texttt{2026-08-21}: the supplied curation labels the problem open and characterizes the extracted prompt as incomplete, concerning Chern classes and cycles supported in a given codimension. It reports failures of a naive formula in surrounding text and no stated replacement theorem settling the extracted problem. This is supplied dated curation, not an independently verified literature assessment.
\item The registry labels the statement \texttt{corrected\_verified} and reports comparison with the official Oberwolfach report. That verification claim is supplied metadata; neither the source nor the mathematical and literature claims have been independently checked for this draft.
\end{itemize}
\end{document}
